\chapter{Définir sa grille en C++}

Le chapitre précédent décrivait un plateau dans un fichier JSON. Ce chapitre
montre l'autre voie — et corrige au passage un malentendu répandu.

\section{Le malentendu : « je suis obligé de passer par du JSON »}

Non. \textbf{Vous ne l'avez jamais été.} Le contrat ne vous impose ni format de
plateau, ni fichier, ni même l'existence d'un fichier. Regardez ce que
\texttt{LoadBoardJson} vaut dans \nkfile{exemples/rules/RegleMinimale.cpp} :

\begin{nkcode}{exemples/rules/RegleMinimale.cpp — V\_LoadBoardJson}
/// Plateau fige : on REFUSE tout chargement. L'atelier l'affiche proprement
/// (« le moteur a REFUSE ce plateau ») au lieu de faire semblant.
int32 V_LoadBoardJson(NkcRules, const char *) { return 0; }
\end{nkcode}

Ce module refuse tout JSON, et il fonctionne parfaitement dans l'atelier. Son
plateau est construit en C++, dans \texttt{Create} :

\begin{nkcode}{exemples/rules/RegleMinimale.cpp — V\_Create}
for (int32 y = 0; y < kSide; ++y)
    for (int32 x = 0; x < kSide; ++x) {
        R->coords[y * kSide + x].q = static_cast<int16>(x);
        R->coords[y * kSide + x].r = static_cast<int16>(y);
    }
\end{nkcode}

\begin{nkretenir}[Trois choses ont TOUJOURS été à vous]
\begin{description}
  \item[La forme du plateau] quelles coordonnées existent : c'est le tableau
        \texttt{coords} que vous exposez dans \texttt{GetView}. Une coordonnée
        absente est du hors-plateau, point.
  \item[Le voisinage] qui touche qui : c'est \emph{votre} code dans
        \texttt{GenerateLegalMoves}. \nkfile{ConquerorGeometry.h} n'est qu'une
        \textbf{commodité}, pas une obligation — vous pouvez ne pas l'inclure.
  \item[Les cases bloquées] c'est \texttt{kCellBlocked} dans le champ
        \texttt{owner} de la vue. Rien d'autre.
\end{description}
Le JSON n'existe que pour qu'un \emph{designer} — quelqu'un qui ne compile pas —
puisse essayer vingt formes dans l'après-midi.
\end{nkretenir}

\section{Ce qui vous échappait vraiment : la projection écran}

Une seule chose ne vous appartenait pas : \textbf{où l'atelier dessine vos
cellules, et quelle forme il leur donne}. Il le déduisait de \texttt{topology},
donc de quatre valeurs possibles — deux hexagones, deux carrés.

Conséquence : un plateau que vos règles savaient parfaitement jouer pouvait être
\textbf{impossible à afficher}. Un plateau circulaire, en triangles, en octogones,
à cellules de tailles inégales : le moteur tournait, l'écran mentait.

L'ABI 3 rend la main au module.

\begin{nkcode}{include/Conqueror/ConquerorRulesABI.h — bloc geometrie (ABI 3)}
// TOUT CE BLOC EST OPTIONNEL : laisser a nullptr fait retomber l'atelier
// sur la projection standard de `topology`.

/// Centre de la cellule `c`, en unites de cellule (1.0 = un « pas » de
/// grille). L'atelier cadre et met a l'echelle ensuite.
int32 (*GetCellCenter)(NkcRules self, NkcCoord c, float32 *outXY) = nullptr;

/// Contour de la cellule `c` : `capPoints` paires (x, y) au plus, en unites
/// de cellule, RELATIVES au centre. Renvoie le nombre de points ecrits.
uint32 (*GetCellShape)(NkcRules self, NkcCoord c, float32 *outXY,
                       uint32 capPoints) = nullptr;
\end{nkcode}

\begin{nkpiege}[Ces flottants ne violent pas « zéro flottant »]
L'exigence \S17.1 porte sur la \textbf{logique de règles} : ce qui décide d'un
coup, alimente un état, entre dans une empreinte. Ces deux fonctions ne font que
de la \textbf{présentation} — l'atelier les appelle pour dessiner, jamais pour
jouer.

Le repère est simple : si votre \texttt{HashState} ne les voit pas, elles sont du
bon côté de la frontière.
\end{nkpiege}

\section{Un plateau circulaire, entièrement en C++}

Le fichier : \nkfile{Applications/ConquerorLab/exemples/rules/GrilleLibre.cpp}.
Trois anneaux concentriques, 19 cellules en secteurs. Ni hexagone, ni carré.

\begin{nkcode}{exemples/rules/GrilleLibre.cpp — la forme, en entiers}
constexpr int32 kRings            = 3;
constexpr int32 kRingSize[kRings] = {1, 6, 12};
constexpr int32 kCells            = 1 + 6 + 12;   // 19
constexpr int32 kRingStart[kRings] = {0, 1, 7};   // index du 1er de chaque anneau

int32 IndexOf(NkcCoord c) {
    if (c.q < 0 || c.q >= kRings) return -1;
    if (c.r < 0 || c.r >= kRingSize[c.q]) return -1;
    return kRingStart[c.q] + c.r;
}
\end{nkcode}

Une coordonnée vaut ici $(q = anneau,\ r = index\ dans\ l'anneau)$. C'est
\emph{notre} convention : le contrat ne dit rien du sens de $q$ et $r$ hors
topologie standard, et c'est exactement ce qui rend la chose possible.

\subsection{Le voisinage, défini par nous}

\begin{nkcode}{exemples/rules/GrilleLibre.cpp — Neighbors}
uint32 Neighbors(NkcCoord c, NkcCoord *out, uint32 cap) {
    uint32 n = 0;
    auto   push = [&](int32 ring, int32 idx) {
          if (ring < 0 || ring >= kRings) return;
          const int32 sz = kRingSize[ring];
          const int32 r  = ((idx % sz) + sz) % sz;   // modulo circulaire
          if (n < cap) { out[n].q = (int16)ring; out[n].r = (int16)r; }
          ++n;
    };

    if (c.q == 0) {                     // le centre touche tout l'anneau 1
        for (int32 i = 0; i < kRingSize[1]; ++i) push(1, i);
        return n;
    }
    push(c.q, c.r - 1);                 // voisin arriere sur l'anneau
    push(c.q, c.r + 1);                 // voisin avant  sur l'anneau
    if (c.q == 1) {
        push(0, 0);                     // vers le centre
        push(2, c.r * 2);               // vers l'exterieur : DEUX cellules
        push(2, c.r * 2 + 1);
    } else {                            // anneau 2
        push(1, c.r / 2);               // vers l'interieur : UNE cellule
    }
    return n;
}
\end{nkcode}

C'est le cœur de la démonstration : \nkfile{ConquerorGeometry.h} n'est même pas
inclus dans ce fichier. Aucune topologie du contrat ne sait exprimer « le centre
touche tout l'anneau 1 », ni « une cellule de l'anneau 1 touche deux cellules de
l'anneau 2 ». Ici, si — et le tout reste \textbf{entièrement entier}, donc le
rejeu bit-à-bit tient.

\begin{nkretenir}
L'ordre reste \textbf{normatif} : \texttt{push} est appelé dans un ordre fixe, et
c'est lui qui fixe l'ordre de génération des coups. Ne le réordonnez pas.
\end{nkretenir}

\subsection{Où dessiner : \texttt{GetCellCenter}}

\begin{nkcode}{exemples/rules/GrilleLibre.cpp — V\_GetCellCenter}
int32 V_GetCellCenter(NkcRules, NkcCoord c, float32 *outXY) {
    if (!outXY) return 0;
    if (c.q < 0 || c.q >= kRings) return 0;
    if (c.q == 0) { outXY[0] = 0.f; outXY[1] = 0.f; return 1; }
    const float32 a = CellAngle(c);
    const float32 R = static_cast<float32>(c.q);   // anneau 1 -> rayon 1
    outXY[0] = R * std::cos(a);
    outXY[1] = R * std::sin(a);
    return 1;
}
\end{nkcode}

Des coordonnées polaires, en \textbf{unités de cellule}. Vous ne vous souciez ni
du zoom, ni des pixels, ni de la taille de la fenêtre : l'atelier calcule la boîte
englobante réelle de vos centres et met tout à l'échelle.

Renvoyer 0 signifie « je ne sais pas placer celle-ci » — l'atelier retombe alors
sur la topologie \textbf{pour cette cellule seulement}. Vous pouvez donc ne
personnaliser qu'une partie du plateau.

\subsection{Quelle forme : \texttt{GetCellShape}}

\begin{nkcode}{exemples/rules/GrilleLibre.cpp — V\_GetCellShape, secteur d'anneau}
// 4 points sur l'arc exterieur, puis 4 sur l'interieur en sens inverse.
for (uint32 i = 0; i < seg; ++i) {
    const float32 t = static_cast<float32>(i) / static_cast<float32>(seg - 1);
    const float32 a = a0 - half + 2.f * half * t;
    outXY[n * 2]     = (R + 0.46f) * std::cos(a) - cx;
    outXY[n * 2 + 1] = (R + 0.46f) * std::sin(a) - cy;
    ++n;
}
for (uint32 i = 0; i < seg; ++i) {
    const float32 t = static_cast<float32>(seg - 1 - i) / static_cast<float32>(seg - 1);
    const float32 a = a0 - half + 2.f * half * t;
    outXY[n * 2]     = (R - 0.46f) * std::cos(a) - cx;
    outXY[n * 2 + 1] = (R - 0.46f) * std::sin(a) - cy;
    ++n;
}
\end{nkcode}

Un polygone \textbf{relatif au centre}, en unités de cellule. Huit points
suffisent à faire un secteur d'anneau convaincant ;
\texttt{kMaxCellPoints} vaut 12.

\begin{nkpiege}[Le polygone doit rester convexe]
L'atelier remplit vos cellules par un éventail de triangles depuis le premier
sommet — c'est la primitive dont dispose \texttt{NkGuiDrawList}, qui n'a pas de
remplissage de polygone quelconque. Un contour concave se remplira de travers.

Un secteur d'anneau assez fin passe très bien ; une cellule en U, non.
\end{nkpiege}

\subsection{Le déclarer aux autres : \texttt{GetNeighbors}}

Votre voisinage est écrit, il marche, vos coups sont bons. Il reste un problème :
\textbf{personne d'autre ne le connaît}.

\begin{itemize}
  \item une IA qui évalue une position (« combien d'ennemis touche cette
        case ? ») n'avait que \texttt{view.topology} pour le deviner. Sur une
        grille libre, elle devinait \textbf{faux, en silence} ;
  \item l'atelier ne pouvait pas montrer le voisinage à l'écran, donc un
        stagiaire qui se trompait d'adjacence n'avait aucun moyen de le
        \textbf{voir}.
\end{itemize}

D'où une troisième entrée — et celle-ci est \textbf{entièrement entière}, c'est
de la règle, pas du dessin :

\begin{nkcode}{exemples/rules/GrilleLibre.cpp — V\_GetNeighbors}
/// Le voisinage, DECLARE. Sans cette entree, une IA qui evalue une position
/// sur ce plateau devinerait l'adjacence a partir de `topology` — et
/// devinerait faux, en silence.
uint32 V_GetNeighbors(NkcRules, NkcCoord c, NkcCoord *out, uint32 cap) {
    return Neighbors(c, out, cap);
}
\end{nkcode}

Une ligne, puisque la fonction existait déjà. Vérifié :

\begin{nkcode}{Sortie du banc de verification, 2026-08-07}
GetNeighbors(0,0) -> 6 voisins : (1,0) (1,1) (1,2) (1,3) (1,4) (1,5)
\end{nkcode}

Le centre touche tout l'anneau 1 — aucune topologie du contrat ne sait exprimer
cela.

\begin{nkretenir}[Le bouton « Voisinage » du panneau Plateau]
Il trace, au survol, les liens de la case vers ses voisins. C'est le \textbf{seul
moyen de vérifier une adjacence à l'œil}, et il lit \texttt{GetNeighbors}. Sans
lui, une erreur d'adjacence se découvre par un coup légal inexplicable, trois
heures plus tard.
\end{nkretenir}

\section{Comment l'atelier s'y adapte}

Tout passe par un objet minuscule, reconstruit à chaque image :

\begin{nkcode}{src/ConquerorLab/NkcBoardRender.h — le projecteur}
struct NkcProjector {
        const NkcRulesVTable *vt       = nullptr;
        NkcRules              inst     = nullptr;
        NkcTopology           topology = NkcTopology::HexPointy;
        bool                  custom   = false;  // le module declare sa geometrie
};
\end{nkcode}

\begin{nkretenir}[Pourquoi il n'a aucun état persistant]
Le rendre persistant obligerait à l'invalider quand le module change — et c'est
exactement le genre d'invalidation qu'on oublie. Le reconstruire coûte deux
affectations et une comparaison de pointeur.
\end{nkretenir}

Une conséquence à connaître, côté picking :

\begin{nkcode}{src/ConquerorLab/NkcBoardRender.h — ProjPickCell}
// Deux chemins, et c'est assume :
//   - geometrie standard : inversion analytique en O(1) (arrondi cube),
//     exacte jusqu'aux aretes ;
//   - geometrie du module : aucune inverse n'existe, on cherche le centre
//     le plus proche parmi les cases. O(n) sur quelques dizaines de cases
//     ne se mesure pas, et c'est la seule methode qui marche pour une
//     grille quelconque.
\end{nkcode}

Le clic sur une grille libre désigne donc la cellule dont le \emph{centre} est le
plus proche, à condition d'être à portée. Sur des cellules très allongées, le
picking peut sembler légèrement décalé près des bords : c'est le prix d'une
géométrie arbitraire, et il est assumé.

\section{JSON ou C++ : lequel choisir}

\begin{center}
\begin{tabular}{@{}p{6.4cm}p{7.4cm}@{}}
\toprule
\textbf{Prenez le JSON quand\ldots} & \textbf{Prenez le C++ quand\ldots} \\
\midrule
la forme se décrit par une liste de coordonnées & la forme se décrit par une
  \emph{formule} (anneaux, spirale, fractale) \\
un designer doit pouvoir l'essayer sans compiler & le voisinage n'est pas
  géométrique (portails, torus, graphe quelconque) \\
vous voulez vingt variantes de la même idée & les cellules n'ont ni la même
  taille ni la même forme \\
\bottomrule
\end{tabular}
\end{center}

Rien n'interdit les deux : votre \texttt{LoadBoardJson} peut lire la liste des
cases, pendant que \texttt{GetCellCenter} en calcule la position. C'est même
souvent le bon partage — la \emph{forme} en donnée, la \emph{projection} en code.

\section{État vérifié}

\begin{nkcode}{Banc d'essai, 2026-08-07 — GrilleLibre contre IAMinimale}
       regles : GrilleLibre (3 anneaux) 1.0.0 (palier 0)
       IA     : IAMinimale 1.0.0
       plateau : 19 cases, 2 joueurs
  OK   l'IA n'a jamais produit de coup illegal
  OK   la partie se termine d'elle-meme        -> 17 coups, vainqueur 0, 10/9
  OK   le journal se rejoue integralement
  OK   rejeu deterministe : empreinte identique -> 15623799979863270834
  OK   « aucun coup legal » equivaut a « joueur bloque »
=== 13 verifications OK, 3 echecs ===
\end{nkcode}

Les trois échecs sont \textbf{attendus} : le banc d'essai vérifie en dur le
plateau du moteur de \emph{référence} (« 42 cases », « 2 totems par joueur », et
la présence du paramètre \texttt{portee\_duplication}). Un plateau de 19 cellules
en anneaux n'a aucune de ces propriétés. Tout ce qui porte sur le \textbf{contrat}
passe.

\begin{nkretenir}[Une leçon de méthode, au passage]
Un test qui code en dur les valeurs d'une implémentation particulière n'est pas un
test du contrat. Si vous reprenez \nkfile{tests/NkcAbiHarness.cpp} pour votre
module, ces trois vérifications-là sont à adapter — les treize autres, non.
\end{nkretenir}

\begin{nkexo}[1 — Le tore]
Partez de \nkfile{RegleMinimale.cpp} et rendez son carré $5\times5$
\emph{torique} : la colonne 4 touche la colonne 0, la ligne 4 touche la ligne 0.
Vous n'avez à toucher que la génération des coups. Le plateau reste affiché à
plat — que manque-t-il pour \emph{voir} le repliement ? Est-ce grave ?
\end{nkexo}

\begin{nkexo}[2 — Le triangle]
Écrivez un plateau en cellules \textbf{triangulaires} : un grand triangle
subdivisé, où chaque petite cellule touche ses trois voisines de côté. Vous aurez
besoin des deux orientations (pointe en haut, pointe en bas) — c'est le vrai
exercice. \texttt{GetCellShape} doit renvoyer trois points.
\end{nkexo}

\begin{nkexo}[3 — La question de fond]
\texttt{GrilleLibre} déclare \texttt{supportsHex = 0} et \texttt{supportsSquare
= 0}. Cherchez dans l'atelier qui lit ces deux champs. Que faudrait-il en faire ?
Proposez une réponse en trois lignes — et vérifiez si le code actuel la respecte.
\end{nkexo}
